% Exercises for Chapter 11 (conditioning and conditional expectation). Statements only:
% solutions are kept privately. Throughout, (Omega, F, P) is a probability space, G is a
% sub-sigma-algebra of F, m is Lebesgue measure, and the notation is that of the notes.

\begin{exercise}[The discrete case]\label{ex:ce-discrete-case}
\easy In 2) of \noteref{rem:ce-density-remarks} the notes say that in the discrete case one can
state and prove a theorem analogous to \noteref{thm:ce-density}. Let $(X,Y)$ be a discrete
random vector with $X\ge0$ or $X\in L^1$, and let $h$ be as in \noteref{def:ce-discrete}, with
$h(y)=0$ where $f_Y(y)=0$. You may use \noteref{rem:ce-discrete-property},
\noteref{thm:ce-doob-dynkin} and \noteref{thm:ce-existence}.
\begin{enumerate}[label=(\alph*)]
  \item Show that $h(Y)=\mathbb E[X\mid\sigma(Y)]$ a.s.: the elementary definition of
    $\mathbb E[X\mid Y]$ and Kolmogorov's agree. What changes in
    \noteref{rem:ce-discrete-property} when $X\ge0$ is not integrable?
  \item Conversely, let $\tilde h$ be Borel with $\tilde h(Y)=\mathbb E[X\mid\sigma(Y)]$ a.s.
    Show that $\tilde h(y)=h(y)$ for every $y$ with $\mathbb P(Y=y)>0$.
  \item Let $E\in\mathcal F$ and $X=\1_E$. Show that
    $\mathbb E[\1_E\mid Y=y]=\mathbb P(E\mid Y=y)=\mathbb P(E\cap\{Y=y\})/\mathbb P(Y=y)$ whenever
    $\mathbb P(Y=y)>0$.
  \item Let $Y$ be a Bernoulli variable, $\mathbb P(Y=0)=\mathbb P(Y=1)=\frac12$, and let the joint
    density of $(X,Y)$ be $f_{(X,Y)}(1,0)=\frac12$ and
    $f_{(X,Y)}(k,1)=\frac1{2k(k+1)}$ for $k\in\mathbb N^*$. Check that this is a joint density,
    compute $h(0)$, $h(1)$ and $\mathbb E[X\mid Y]$, and verify property 2) of
    \noteref{def:ce-conditional-expectation} on each of the four events of $\sigma(Y)$. Is
    $X\in L^1$?
\end{enumerate}
\end{exercise}

\begin{exercise}[Beyond integrable variables]\label{ex:ce-beyond-L1}
\hard This exercise completes the remarks of \noteref{def:ce-conditional-expectation} and of
the proof of \noteref{thm:ce-existence}. For $X\ge0$ let $Q(A)=\int_AX\,\dd\mathbb P$ for
$A\in\mathcal G$.
\begin{enumerate}[label=(\alph*)]
  \item Let $X\ge0$. Show that $Q$ is $\sigma$-finite on $(\Omega,\mathcal G)$ if and only if
    $\mathbb E[X\mid\mathcal G]<+\infty$ a.s.
  \item Let $X$ be a random variable with $X^-\in L^1$ (the notes say that this suffices).
    Show that $\hat X=\mathbb E[X^+\mid\mathcal G]-\mathbb E[X^-\mid\mathcal G]$ is defined a.s.,
    with values in $(-\infty,+\infty]$, that it satisfies 1) and 2) of
    \noteref{def:ce-conditional-expectation} (the integrals in 2) taking values in
    $(-\infty,+\infty]$), and that any $\mathcal G$-measurable $\tilde X:\Omega\to(-\infty,+\infty]$
    satisfying 2) equals $\hat X$ a.s. \emph{Hint:} first show that $\tilde X^-\in L^1$, then
    compare $\hat X$ and $\tilde X$ on the sets $\{\hat X<\tilde X\le n\}$ and
    $\{\hat X\le n,\ \tilde X=+\infty\}$.
  \item Let $Y$ be a Bernoulli variable with $\mathbb P(Y=1)=\frac12$, let $W\ge0$ be independent of
    $Y$ with $\mathbb E[W]=+\infty$, and let $X=YW-1$. Show that $X^-\in L^1$ and
    $X\notin L^1$, and that $\mathbb E[X\mid Y]=+\infty$ on $\{Y=1\}$ and $-1$ on $\{Y=0\}$. You
    may use \noteref{thm:pm-expectation-product} for non-negative variables, in the form
    $\mathbb E[UV]=\mathbb E[U]\,\mathbb E[V]$ in $[0,+\infty]$ for independent $U,V\ge0$.
\end{enumerate}
\end{exercise}

\begin{exercise}[A finite partition]\label{ex:ce-partition}
\medium This is partly a computational exercise. Let $A_1,\dots,A_k\in\mathcal F$ be a
partition of $\Omega$ into non-empty sets, and $\mathcal G=\sigma(A_1,\dots,A_k)$.
\begin{enumerate}[label=(\alph*)]
  \item Show that $\mathcal G$ consists of the $2^k$ unions of cells, and that a function
    $\Omega\to\overline{\mathbb R}$ is $\mathcal G$-measurable if and only if it is constant on
    each $A_i$. How does this relate to \noteref{thm:ce-doob-dynkin}, with
    $Y=\sum_ii\1_{A_i}$?
  \item Let $X\ge0$ or $X\in L^1$. Show that
    \[
      \mathbb E[X\mid\mathcal G]=\sum_{i:\,\mathbb P(A_i)>0}\frac{\mathbb E[X\1_{A_i}]}{\mathbb P(A_i)}\,\1_{A_i}
      \qquad\text{a.s.},
    \]
    with any values on the cells of probability $0$.
  \item Two fair dice are rolled: $\Omega=\{1,\dots,6\}^2$ with the uniform probability, $D_1,D_2$
    the two results, $M=\max(D_1,D_2)$ and $S=D_1+D_2$. In exact rational arithmetic, write a
    function that computes $\mathbb E[D_1\mid\mathcal G]$ for $\mathcal G$ generated by a
    partition, and one that checks property 2) of \noteref{def:ce-conditional-expectation} on
    \emph{every} event of $\mathcal G$. Use them to show that
    $\mathbb E[D_1\mid M=m]=\frac{m(3m-1)}{2(2m-1)}$ and $\mathbb E[D_1\mid S]=\frac S2$, and prove
    both formulas by hand.
  \item Show that $\frac S2$ satisfies property 2) on all the $64$ events of $\sigma(M)$, yet it is
    not $\mathbb E[D_1\mid M]$. Which property fails? Explain the first fact by a symmetry.
  \item Compute the variance of $\mathbb E[D_1\mid\mathcal G]$ for $\mathcal G=\{\emptyset,\Omega\}$,
    $\sigma(M)$, $\sigma(S)$, $\sigma(M,\min(D_1,D_2))$ and $\mathcal F$, and order them. Two of
    the five conditional expectations coincide although the $\sigma$-algebras differ: which, and
    why?
\end{enumerate}
\end{exercise}

\begin{exercise}[Conditioning on a sum]\label{ex:ce-sums}
\medium (in the style of Ross, Ch.~6) \noindent
\begin{enumerate}[label=(\alph*)]
  \item Let $X\sim B(n_1,p)$ and $Y\sim B(n_2,p)$ be independent, $0<p<1$, and $Z=X+Y$, so that
    $Z\sim B(n_1+n_2,p)$ (\noteref{prop:rv-notable-discrete}). Show that for
    $0\le m\le n_1+n_2$
    \[
      f_{X\mid Z=m}(j)=\frac{\binom{n_1}j\binom{n_2}{m-j}}{\binom{n_1+n_2}m}
      \qquad\text{for }\max(0,m-n_2)\le j\le\min(n_1,m),
    \]
    which does not depend on $p$, and interpret it with an urn. Deduce
    $\mathbb E[X\mid Z]=\frac{n_1}{n_1+n_2}Z$.
  \item Let $X,Y$ be independent with the geometric law $\mathbb P(X=k)=p(1-p)^{k-1}$,
    $k\in\mathbb N^*$. Show that, given $X+Y=m$ ($m\ge2$), $X$ is uniform on
    $\{1,\dots,m-1\}$, and find $\mathbb E[X\mid X+Y]$. Compare with
    \noteref{exa:ce-exponential}.
  \item Let $X,Y\in L^1$ be independent with the same law, and $Z=X+Y$. Show, directly from
    \noteref{def:ce-conditional-expectation} and without densities, that
    $\mathbb E[X\mid Z]=\frac Z2$ a.s. \emph{Hint:} by \noteref{thm:pm-independence-product},
    $(X,Y)$ and $(Y,X)$ have the same law; deduce that so do $(X,Z)$ and $(Y,Z)$.
  \item Let $X,Y$ be independent and uniform on $(0,1)$. Find the density of $X$ given $Z=z$ for
    $0<z<2$ by \noteref{thm:ce-density}, and check (c).
\end{enumerate}
\end{exercise}

\begin{exercise}[How many uniforms to exceed $x$]\label{ex:ce-renewal}
\hard (in the style of Ross, Ch.~7) Let $U_1,U_2,\dots$ be independent and uniform on $(0,1)$,
$S_n=U_1+\dots+U_n$, and for $0\le x\le1$
\[
  N(x)=\min\{n\ge1:S_n>x\},\qquad m(x)=\mathbb E[N(x)] .
\]
\begin{enumerate}[label=(\alph*)]
  \item Show by induction, with Tonelli's theorem, that $\mathbb P(S_n\le x)=\frac{x^n}{n!}$ for
    $0\le x\le1$. Deduce that $m(x)=e^x$ and $\mathbb E[N(1)^2]=3e$. \emph{Hint:}
    $\mathbb E[N]=\sum_{n\ge0}\mathbb P(N>n)$ for $N$ with values in $\mathbb N^*$.
  \item Now compute $m$ by conditioning on $U_1$. Let
    $N'(t)=1+\sum_{n\ge1}\1_{\{U_2+\dots+U_{n+1}\le t\}}$ for $0\le t\le1$. Show that
    $N(x)=1+N'(x-U_1)\1_{\{U_1\le x\}}$, and that
    \[
      \mathbb E[N(x)\mid U_1]=h(U_1),\qquad h(u)=\begin{cases}1+m(x-u)&\text{if }u\le x,\\1&\text{if }u>x,\end{cases}
    \]
    by checking property 2) of \noteref{def:ce-conditional-expectation} term by term in the sum.
    (Why does $\mathbb E[N'(t)]=m(t)$?)
  \item Deduce that $m(x)=1+\int_0^xm(t)\,\dd t$, and solve this equation.
  \item Show that $\operatorname{Var}(N(1))=3e-e^2\approx0.7658$, and find the number $n$ of
    simulations for which the asymptotic $95\%$ confidence interval for $\mathbb E[N(1)]$ has
    half-width at most $0.002$. Run them (seed $2026$): does the interval contain $e$?
\end{enumerate}
\end{exercise}

\begin{exercise}[A triangular density]\label{ex:ce-joint-density}
\medium (in the style of Ross, Ch.~6) This is partly a computational exercise. Let $(X,Y)$ have
the joint density $f(x,y)=8xy$ for $0<x<y<1$, and $0$ elsewhere.
\begin{enumerate}[label=(\alph*)]
  \item Find $f_X$, $f_Y$, $f_{X\mid Y=y}$ and $f_{Y\mid X=x}$, and show that
    \[
      \mathbb E[X\mid Y]=\frac23Y,\qquad\mathbb E[Y\mid X]=\frac23\,\frac{1+X+X^2}{1+X}.
    \]
    Are $X$ and $Y$ independent (5) of \noteref{rem:ce-density-remarks})? Compute
    $\mathbb E[X]$ directly and from $\mathbb E[X\mid Y]$, using 2) of
    \noteref{def:ce-conditional-expectation} with $A=\Omega$.
  \item Compute $\mathbb P_{X\mid Y=1/2}\big((-\infty,\frac14]\big)$ and
    $\mathbb P_{Y\mid X=1/2}\big([\frac34,+\infty)\big)$.
  \item Let $U,V$ be independent and uniform on $(0,1)$, $Y=U^{1/4}$ and $X=Y\sqrt V$. Show
    that $(X,Y)$ has density $f$. (This is \noteref{eq:ce-joint-from-conditional} read backwards:
    first draw $Y$ from $f_Y$, then $X$ from $f_{X\mid Y=y}$.)
  \item Simulate $10^6$ pairs (seed $2026$). Compare the empirical distribution function of $X$
    with $F_X(x)=2x^2-x^4$, with a tolerance from the Dvoretzky--Kiefer--Wolfowitz inequality
    at level $99\%$. Compare the averages of $Y$ over the bins $[\frac{k}{20},\frac{k+1}{20})$
    of $X$ with the exact values $\mathbb E[Y\1_{\{X\in I\}}]/\mathbb P(X\in I)$, in units of
    standard errors.
\end{enumerate}
\end{exercise}

\begin{exercise}[Conditioning by binning]\label{ex:ce-binning}
\medium This is a computational exercise. Let $X,Y$ be independent $\mathrm{Exp}(1)$ and $Z=X+Y$,
so that $\mathbb E[X\mid Z]=\frac Z2$ (\noteref{exa:ce-exponential}). Let
$I_k=[0.2(k-1),0.2k)$, $k=1,\dots,15$.
\begin{enumerate}[label=(\alph*)]
  \item Show that $\mathbb E[X\1_{\{Z\in I\}}]/\mathbb P(Z\in I)=\frac12\mathbb E[Z\1_{\{Z\in I\}}]/\mathbb P(Z\in I)$
    for every interval $I$ with $\mathbb P(Z\in I)>0$. Compute it for $I=I_1$ and compare with
    $\frac12$ of the centre of $I_1$.
  \item Given $Z\in I_k$, $X$ takes values in $[0,0.2k)$, so its conditional variance is at most
    $(0.2k)^2/4$. Use this bound to choose $n$ so that, if bin $I_k$ receives
    $n\,\mathbb P(Z\in I_k)$ samples, the standard error of the average of $X$ over every bin is
    at most $0.02$. Which bin decides $n$?
  \item Simulate $n$ pairs (seed $2026$) and compute the bin averages of $X$, their standard
    errors, and the largest deviation from the exact values of (a), in standard errors. Compare
    the averages over the first three bins with half the bin centres. Is the bound of (b)
    sharp?
  \item Estimate $\mathbb E[(X-h(Z))\1_{\{Z\in B\}}]$ for $B=[0,1)$, $[1,2)$, $[2,+\infty)$, with
    $h(z)=\frac z2$ and with $h(z)=\frac z3$, and compare with the exact values.
  \item Let now $X\sim\Gamma(2,1)$ and $Y\sim\mathrm{Exp}(1)$ be independent. Show with
    \noteref{prop:rv-gamma-beta} that $\mathbb E[X\mid Z]=\frac23Z$, so that the exact bin values are
    $\frac23\mathbb E[Z\1_{\{Z\in I_k\}}]/\mathbb P(Z\in I_k)$ with $Z\sim\Gamma(3,1)$. With the same $n$
    and bins, compute the bin averages of $X$, their standard errors, and the largest deviation
    from the exact values. Is every standard error still at most $0.02$, although $n$ was chosen
    for the exponential pair?
\end{enumerate}
\end{exercise}

\begin{exercise}[A Poisson pair, simulated]\label{ex:ce-poisson-simulation}
\easy This is a computational exercise on \noteref{exa:ce-poisson}, with $\lambda=2$ and
$\mu=3$. Simulate $10^6$ pairs $(X,Z)$ (seed $2026$).
\begin{enumerate}[label=(\alph*)]
  \item Compare the relative frequencies of $X=k$ among the pairs with $Z=5$ with the law
    $B(5,0.4)$, in units of standard errors.
  \item For $m=0,\dots,12$, compare the average of $X$ over the pairs with $Z=m$ with
    $h(m)=0.4\,m$. How many deviations exceed two standard errors, and how many would you expect?
  \item Compute exactly $\mathbb E[X\1_{\{Z\le4\}}]$, using property 2) of
    \noteref{def:ce-conditional-expectation}, and compare with its Monte Carlo estimate.
  \item Show that $Z-k$, given $X=k$, has the law $\mathcal P(\mu)$, so that
    $\mathbb E[Z\mid X]=X+\mu$; check it on the simulation. Why is $\mathbb E[Z\mid X]$ not
    $\frac{\lambda+\mu}\lambda X$, although $\mathbb E[X\mid Z]=\frac\lambda{\lambda+\mu}Z$?
\end{enumerate}
\end{exercise}

\begin{exercise}[Functions of $Y$]\label{ex:ce-doob-dynkin}
\medium This exercise extends \noteref{thm:ce-doob-dynkin}; the last part is computational.
\begin{enumerate}[label=(\alph*)]
  \item Show that the theorem holds, with the same proof, for $Y:\Omega\to\mathbb R^d$ (with
    $h:\mathbb R^d\to\mathbb R$) and for $X:\Omega\to[0,+\infty]$ (with
    $h:\mathbb R\to[0,+\infty]$). Which steps change?
  \item Show that if $h_1(Y)=h_2(Y)$ everywhere, then $h_1=h_2$ on $Y(\Omega)$, and that if
    $h_1(Y)=h_2(Y)$ a.s., then $h_1=h_2$ $\mathbb P_Y$-a.e. Give an example with $h_1(Y)=h_2(Y)$
    everywhere but $h_1\ne h_2$ on a set of positive Lebesgue measure.
  \item Let $Y$ be symmetric (that is, $Y$ and $-Y$ have the same law) with $\mathbb P(Y\ne0)>0$.
    Show that $Y$ is not $\sigma(Y^2)$-measurable, and that $\mathbb E[Y\mid Y^2]=0$ if
    $Y\in L^1$.
  \item On the space of two dice of \exref{ex:ce-partition}, write a function that, given $Z$ and
    $Y$ on a finite space, returns $h$ with $Z=h(Y)$ or a level set of $Y$ on which $Z$ is not
    constant. Decide whether $D_1+D_2$ is a function of $D_1D_2$; whether $\max(D_1,D_2)$ is a
    function of $(D_1+D_2,|D_1-D_2|)$; and whether $D_1$ is a function of
    $(D_1+D_2,|D_1-D_2|)$. Give $h$ when it exists.
\end{enumerate}
\end{exercise}
